\documentclass[10pt,a4paper]{amsart}
\usepackage[margin=19mm]{geometry}
\usepackage{amsmath,amssymb,mathtools}
\usepackage[scr=rsfs,cal=euler]{mathalfa}
\usepackage{xfrac}
\usepackage[unicode,hidelinks,bookmarks=false]{hyperref}
\newcommand{\ie}{i.e.\ }
\newcommand{\cf}{cf.\ }
\newcommand{\resp}{respectively\ }
\newcommand{\Subsection}{\subsection}
\providecommand{\nobreakdash}{\nobreak\hbox{-}}
\setlength{\emergencystretch}{2em}
\multlinegap0pt
\usepackage{breqn}
\newtheorem{prop}{Proposition}
\newtheorem{thm}{Theorem}
\usepackage{cleveref}
\crefname{prop}{Proposition}{Propositions}
\crefname{thm}{Theorem}{Theorems}
\newcommand{\C}{\mathbb{C}}
\def\MV{{\rm MV}}
\newcommand{\R}{\mathbb{R}}
\newcommand{\Z}{\mathbb{Z}}
\newcommand{\abs}[1]{\left| #1 \right|}
\newcommand{\be}{\begin{equation}}
\newcommand{\ee}{\end{equation}}
\renewcommand{\hat}{\widehat}
\DeclareMathOperator{\proj}{proj}
\title{Preserving Hodge vectors of lattice polytopes}
\author{Vadym Kurylenko, Benjamin Nill}
\date{Source: arXiv:2602.20765v2 (CC BY 4.0)}
\begin{document}
\maketitle

\noindent\emph{Source and licence.} Preserving Hodge vectors of lattice polytopes. Version arXiv:2602.20765v2 (CC BY 4.0), distributed by arXiv under the Creative Commons Attribution 4.0 International licence, http://creativecommons.org/licenses/by/4.0/, as recorded in the arXiv licence metadata for this exact version and shown on the abstract page https://arxiv.org/abs/2602.20765v2. Full licence notice: \texttt{licenses/arxiv-2602.20765-CC-BY-4.0.txt}.\par\smallskip

\noindent\textit{Editorial scope.} Counted unit: the article's thm thm:main-local (source0.tex lines 334--355). The article's complete original proof of it is delivered with it (source0.tex lines 357--396, one \texttt{proof} environment of 40 lines). No mathematics is written, repaired or re-derived: the statement and the proof are the article's own lines, in the article's own notation, and no retained line is substituted or re-flowed.  Retained uncounted as the source-local context the proof needs: lines 190--192; lines 270--272.  Nothing was omitted from any retained span, so the package carries no omission record.  No retained line contains a citation, so the wrapper carries no bibliography and no bibliography entry.  No retained line contains a figure environment, an image-inclusion command or drawing code, so no image is delivered and the package declares no asset dependency.  Editorial formatting only, in addition: an AMS \texttt{amsart} preamble that restates the article's own declarations and macros used by the retained lines, namely the environments \texttt{prop}, \texttt{thm} on separate counters, together with the standard packages those lines need.  The retained environments print in this document as Proposition 1, Theorem 1 in order of appearance, whereas the article prints the same items with the numbering of its own full document.  The complete native source of the article as distributed in the arXiv e-print for this exact version is bundled unchanged as \texttt{sources/195214/source0.tex}.
\begin{equation}  E(Z_P; u,v) = \frac{1}{uv} \left[(-1)^{d-1} h^*(P; u,v) + (uv- 1)^{d} \right].
\label{E-hyper}
\end{equation}

 \begin{prop} \label{prop: mixed_h_relation} In this situation, the following holds:
   \[ E (Y; u,v) =\frac{1}{(uv)^k} \sum_{I \subseteq[k]} (-1)^{d_I -\abs{I}} (uv-1)^{d-d_I} h^*(P^I; u,v).\]  
\end{prop}

\begin{thm} \label{thm:main-local} \label{Cayley-formula}
Let $P_1,\ldots,P_k, P_{k+1} \subset \R^d$ be lattice polytopes 
such that \[\dim(P_1 + \cdots + P_k) \le k \text{ and } \dim(P_1 + \cdots + P_{k+1})=d \geq 2.\]
Let $U$ be any $k$-dimensional rational subspace of $\R^d$ containing $P_1, \ldots, P_k$ if $k \le d$, and set $U := \R^d$ otherwise.
We denote by $V$ the mixed volume $\MV(P_1, \ldots, P_k)$ of $P_1, \ldots, P_k$ if $k \le d$, and set $V := 0$ otherwise.
 Then the bivariate $h^*$-polynomial of the Cayley polytope $P_1 * \cdots *  P_{k+1}$ equals
\be V (uv)^k h^*(\proj_U(P_{k+1}); u,v) + \sum_{I \subsetneq[k+1]: k+1 \in I} (-1)^{k -\abs{I}} (1-uv)^{d-d_I} h^*(P^I; u,v).
\label{main-eq}
\ee
In particular, we have
\be \ell^*(P_1 * \cdots *  P_{k+1};u,v) = V (uv)^k \ell^*(\proj_U(P_{k+1});u,v)
\label{main-local}\ee
in each of the following two situations: \begin{enumerate}
    \item For each $I \subsetneq[k+1]$ with $k+1 \in I$ we have 
    \begin{itemize}
        \item $|I|+(d-d_I) \le k$ or
        \item $|I|+(d-d_I)=k+1$ and $P^I$ is thin. 
    \end{itemize}
    For instance, this holds if $\dim(P_{k+1})=d$.
    \item $k\ge d$, which implies that $P_1 * \cdots *  P_{k+1}$ is trivially thin.
    \end{enumerate}
\end{thm}

\begin{proof}
Let us consider the main case $k\leq d$ first. 
Let $f_1, \ldots, f_{k+1} \in \C[x^\pm_1, \ldots, x^\pm_d]$ be generic polynomials with Newton polytopes $P_1,\ldots, P_{k+1}$. As $U$ is a rational subspace, we can choose an appropriate lattice transformation of $\Z^d$, so we can assume that $U$ is simply the subspace generated by the first $k$ standard basis vectors and $f_1, \ldots, f_k \in \C[x^\pm_1, \ldots, x^\pm_k]$.


We denote the solution set of $f_1, \ldots, f_k$ in $(\C^*)^k$ by $\hat{Y}$. Note that $\hat{Y}$ is a finite set. It follows from the BKK-theorem that its size equals the mixed volume $V := \MV(P_1, \ldots, P_k)$. In particular, applying \cref{prop: mixed_h_relation} to $P_1, \ldots, P_k$ and using the notation of \cref{prop: mixed_h_relation} yields
   \begin{equation}V =\frac{1}{(uv)^k} \sum_{I \subseteq[k]} (-1)^{d_I -\abs{I}} (uv-1)^{k-d_I} h^*(P^I; u,v).\label{mv}\end{equation}   
Applying \cref{prop: mixed_h_relation} to $P_1, \ldots, P_{k+1}$ gives
   \[ E (Y; u,v) =\frac{1}{(uv)^{k+1}} \sum_{I \subseteq[k+1]} (-1)^{d_I -\abs{I}} (uv-1)^{d-d_I} h^*(P^I; u,v).\]
By taking \eqref{mv} into account, we get that $E(Y; u,v)$ equals
\begin{equation}
\frac{1}{(uv)^{k+1}} \left((uv-1)^{d-k} (uv)^k V + \sum_{I \subseteq[k+1]: k+1 \in I} (-1)^{d -\abs{I}} (1-uv)^{d-d_I} h^*(P^I; u,v)\right).
\label{E-eq1}
\end{equation}


On the other hand, $Y$ is stratified in $\{x'\} \times Z_{x'}$ for $x' \in \hat{Y}$, where $Z_{x'}$ is the hypersurface in $(\C^*)^{d-k}$ defined by the polynomial $g_{x'}(z):=f_{k+1}(x',z) \in \C[z^\pm_1, \ldots, z^\pm_{d-k}]$. Note that each $g_{x'}$ is a generic\footnote{
These Laurent polynomials are nondegenerate with respect to  $\text{Newt}(g_{x'}(z))$ for each $x'$, provided that the coefficients of $f_{k+1}$ are chosen appropriately. Specifically, let $f_{k+1}= \sum_{m \in P_{k+1} \cap \Z^d} c_m x^m$. For $f_{k+1}$ to be nondegenerate with respect to $P_{k+1}$, the coefficient vector $(c_m)_{m \in P_{k+1}}$ must lie in the complement of a finite collection of hypersurfaces in the coefficient space $\mathbb{C}^{|P_{k+1} \cap \mathbb{Z}^d|}$. Similarly, for each $x'$, the requirement that $g_{x'}(z)$ be nondegenerate with respect to its own Newton polytope excludes another finite set of hypersurfaces. Consequently, there exists a Zariski open subset of the coefficient space in which $f_{k+1}$ and all $g_{x'}(z)$ are simultaneously nondegenerate. } polynomial with respect to its Newton polytope which equals the projection $\proj_U(P_{k+1})$ of $P_{k+1}$ along $U$.  We may assume $k<d$ as otherwise (in case (2)), $Z_{x'}$ and hence $Y$ is empty. Note that $\proj_U(P_{k+1})$ is a full-dimensional lattice polytope in $\R^{d-k}$. By \eqref{E-hyper} each $Z_{x'}$ has the same Hodge-Deligne polynomial and as $|\hat{Y}|=V$ we get
\begin{equation}
E (Y; u,v) =  \frac{V}{uv} \left[(-1)^{d-k-1} h^*(\proj_U(P_{k+1}); u,v) + (uv- 1)^{d-k} \right].
\label{E-eq2}
\end{equation}

Equating \eqref{E-eq1} and \eqref{E-eq2} and isolating $h^*(P_1 * \cdots *  P_{k+1};u,v)$ we get our desired formula after some cancellation. 

For the case when $k>d$, the equations  \eqref{mv} and  \eqref{E-eq1} that one derives 
from \cref{prop: mixed_h_relation} still hold with $V=0$. Thus, from these two equations one arrives  at 
\[ h^*(P_1 * \ldots * P_{k+1};u,v) =\sum_{I \subsetneq[k+1]: k+1 \in I} (-1)^{k -\abs{I}} (1-uv)^{d-d_I} h^*(P^I; u,v). \] 



For the additional statements we recall that the Cayley polytopes $P^I$ have dimensions $d_I+|I|-1$, $\proj_U(P_{k+1})$ has dimension $d-k$, and the local $h^*$-polynomial of an $n$-dimensional lattice polytope equals the top degree part of its bivariate $h^*$-polynomial which has degree $n+1$. So, we need to consider the degree $d+k+1$ part in \eqref{main-eq}.
The first term in \eqref{main-eq} has exactly $d+k+1$ as maximal degree and $ V (u v)^k \ell^*(\proj_U(P_{k+1};u,v)$ as the top degree part. Therefore, to arrive at \eqref{main-local} we need that for each $I$ the parts of $h^*(P^I;u,v)$  with degrees $d+k+1-2j$ for $j=0,\ldots, d-d_I$ vanish. 
The easiest possible case is when the maximal degree of $h^*(P^I;u,v)$, i.e., $d_I+\abs{I}$, is smaller than $2d_I-d+k+1$, i.e., $|I|+(d-d_I) \le k$. 
Another notable situation, is when the maximal degree of $h^*(P^I;u,v)$ is exactly  $2d_I-d+k+1$, but the part of $h^*(P^I;u,v)$ of this degree vanishes, i.e. $P^I$ is thin.  

In the situation (2) note that the degree of the Cayley polytope $P_1 * \cdots *  P_{k+1}$ is at most $d$ while its dimension is $d+k\ge 2d$, so it is trivially thin. The projection $\proj_U(P_{k+1})$ is just a point, so its Hodge vector is also $0$. 


\end{proof}

\end{document}
